\section{Overview}
\label{sec:overview}

The Client Interface (CIF) accepts AXI read and write transactions, validates
them, divides each accepted transaction into up to 64 MC request packets (one
per 64-B AXI data beat), and returns AXI responses in order for each AXID.
Read and write traffic use independent instances of the same
transaction-tracking ROB architecture.

\subsection{Global assumptions}
\begin{itemize}[noitemsep]
  \item AXID is 6 bits (64 AXI IDs); each ROB instance has 64 per-AXID Slot
        FIFOs, each of depth four.
  \item One MC request packet carries one 512-bit / 64-B AXI data beat.  One
        accepted AXI transaction produces up to 64 MC packets --- a beat
        \emph{is} an MC packet.
  \item Every AXI data beat transfers a full 64-B unit (\texttt{AxSIZE = 6})
        at a 64-B-aligned address.  Like the other AXI-protocol constraints,
        this is a design-time assumption verified by assertion, not a
        runtime check (see the Request Validator); it is what makes one AXI
        beat correspond to exactly one 64-B MC packet and
        \texttt{exp\_pkts = AxLEN + 1}.
  \item An AXI transaction must not cross 4~KB, so it produces at most
        \texttt{MAX\_PKTS = 4\,KB / 64\,B = 64} MC packets.  With
        runtime-configured \texttt{row\_size} of at least 1~KB, row
        crossings are at most three and affect each packet's address, not
        the packet count.
  \item MC Core echoes the packet's \texttt{GLOBAL\_TAG} on completion, one
        response per accepted 64-B request packet (\S\ref{sec:cifmc-req}).
        Completion \texttt{status} encoding is frozen: \texttt{0}/\texttt{1}/\texttt{2}
        $=$ OK/SLVERR/DECERR (\S\ref{sec:cifmc}).
\end{itemize}

\section{Parameters}
\begin{center}
\begin{tabular}{ll}
\toprule
\textbf{Parameter} & \textbf{Value}\\
\midrule
\texttt{AXID\_W} & 6\\
\texttt{ROB\_DEPTH} & 256 transactions per ROB instance\\
\texttt{ROB\_INDEX\_W} & 8\\
\texttt{DATA\_W} & 512: one 64-B AXI data beat per MC packet\\
\texttt{PACKET\_NUMBER\_W} & 6\\
\texttt{GLOBAL\_TAG\_W} & 14: \texttt{\{ROB\_INDEX[7:0], PACKET\_NUMBER[5:0]\}}\\
\texttt{COMPLETION\_BITMAP\_W} & 64\\
\texttt{MAX\_PKTS} & 64\\
Per-AXID Slot FIFO & 64 FIFOs $\times$ depth 4, per ROB instance\\
Read Data Buffer watermark & 128 entries, owned independently by that buffer\\
\bottomrule
\end{tabular}
\end{center}

\section{Transaction identity and packet formation}
\label{sec:identity}

At AXI transaction acceptance, the applicable ROB Free List allocates
\texttt{ROB\_INDEX[7:0]}.  The index is pushed into the accepting AXID's Slot
FIFO only after the Burst Splitter has supplied \texttt{exp\_pkts}, the
number of 64-B MC packets for the transaction (\texttt{1..64}).  The ROB
initializes \texttt{exp\_pkts} and clears its 64-bit completion bitmap
before the entry becomes eligible for retirement.

At packet emission, the Burst Splitter assigns \texttt{PACKET\_NUMBER}
(\texttt{0..exp\_pkts-1}) and Packet Formation attaches
\texttt{GLOBAL\_TAG = \{ROB\_INDEX[7:0], PACKET\_NUMBER[5:0]\}} (14 bits).
This is the complete MC completion identity.

\section{Write path}
\label{sec:writepath}
\subsection{Request path}

The AW Metadata FIFO, W Data FIFO, Request Validator, Error Handler, Burst
Splitter, Packet Formation Unit, and Write Packet FIFO retain their existing
roles.  The accepting write ROB allocates the transaction index.  The splitter
generates up to 64 MC packets, one per 64-B AXI data beat, and the Write
Packet FIFO forwards them to MC Core.

\subsection{Write ROB and response}

The write ROB has its own 256-entry Free List, Metadata RAM, and 64 Slot
FIFOs.  Its Metadata RAM is directly indexed by \texttt{ROB\_INDEX} and
stores \texttt{exp\_pkts[6:0]}, \texttt{completion\_bitmap[63:0]}, aggregate
\texttt{resp\_status[1:0]}, and the Issue-\#11 footprint
\texttt{wr\_line\_base[25:0]}/\texttt{wr\_line\_count[6:0]}
(\S\ref{sec:samelinehaz}).  MC completions set the bit selected by
\texttt{PACKET\_NUMBER} and update aggregate status.  The \texttt{status}
encoding (\texttt{0}/\texttt{1}/\texttt{2} $=$ OK/SLVERR/DECERR, codes
\texttt{3..15} reserved $\rightarrow$ DECERR) and the worst-status precedence
(\texttt{DECERR} $>$ \texttt{SLVERR} $>$ \texttt{OKAY}) are frozen in
\S\ref{sec:cifmc}.

When the head transaction of an AXID Slot FIFO has \texttt{completion\_bitmap}
bits \texttt{[exp\_pkts-1:0]} all set, retirement selects it in AXID
order.  It holds BRESP valid until the
downstream AXI response path accepts it.  Only that acceptance pops the Slot
FIFO and returns \texttt{ROB\_INDEX} to the write Free List.  Write traffic
has no Read Data Buffer.

\section{Read path}
\label{sec:readpath}
\subsection{Request path}

The AR Metadata FIFO, Request Validator, Error Handler, Burst Splitter,
Packet Formation Unit, and Read Packet FIFO retain their existing roles.  The
read ROB allocates \texttt{ROB\_INDEX} at acceptance, subject to the same-line
read/write hazard check at the AR Metadata FIFO head
(\S\ref{sec:samelinehaz}).  The splitter emits up to 64 MC packets, each
carrying a distinct \texttt{PACKET\_NUMBER} (\texttt{0..exp\_pkts-1}).

\subsection{Read ROB, Data Buffer, and serializer}

The independent read ROB has the same generic structures as the write ROB.
Its Metadata RAM additionally stores \texttt{rd\_dbuf\_base}, the base of a
single contiguously-allocated RD\_SRAM range of \texttt{exp\_pkts} 64-B
lines, set at allocation.  On completion, packet \texttt{PACKET\_NUMBER}'s
returned 64-B data unit is written to \texttt{rd\_dbuf\_base + PACKET\_NUMBER}
and \texttt{completion\_bitmap[PACKET\_NUMBER]} is set.

Once a Slot FIFO head is complete, a Retiring Register holds its
\texttt{ROB\_INDEX} while the Read Serializer emits
\texttt{rd\_dbuf\_base + 0 .. rd\_dbuf\_base + exp\_pkts-1} in packet order,
one AXI beat per 64-B MC packet.  \texttt{RLAST} is asserted on the final
beat (\texttt{PACKET\_NUMBER == exp\_pkts - 1}); only after that beat is
accepted and the \texttt{rd\_dbuf\_base} range released may the Slot FIFO be
popped and \texttt{ROB\_INDEX} returned to the read Free List.  RD\_SRAM
depth (hence \texttt{RD\_SRAM\_ADDR\_W}), its handshakes and latency remain
\textbf{TBD}; its 128-entry watermark is its own capacity control, not ROB
occupancy control.

\subsection{Same-line read/write ordering}
\label{sec:samelinehaz}

The MC core may reorder a read ahead of an older write to the same 64\,B
line.  CIF prevents the resulting stale read by holding a conflicting read at
admission until the overlapping write has retired.  The check is request-wise
and uses only the raw 64\,B line address; no hashed \texttt{daddr} and no CAM
are introduced.

\textbf{Write footprint.} Each outstanding write is tracked in its
\texttt{ROB\_WR} entry as one contiguous 64\,B line interval
\texttt{\{wr\_line\_base[25:0], wr\_line\_count[6:0]\}}, derived from
\texttt{\{AWADDR, AWLEN, AWSIZE, AWBURST\}}:
\begin{itemize}[noitemsep]
  \item \texttt{nbytes = (AWLEN + 1) << AWSIZE}
  \item \texttt{INCR}:  \texttt{start = AWADDR};
        \texttt{last = AWADDR + nbytes - 1}
  \item \texttt{WRAP}:  \texttt{start = AWADDR \& \textasciitilde(nbytes - 1)};
        \texttt{last = start + nbytes - 1}
  \item \texttt{FIXED}: \texttt{start = AWADDR};
        \texttt{last = AWADDR + (1 << AWSIZE) - 1}
  \item \texttt{wr\_line\_base  = start[31:6]}
  \item \texttt{wr\_line\_count = last[31:6] - start[31:6] + 1}, range
        \texttt{[1,64]}
\end{itemize}
Rounding outward to line granularity is mandatory: burst alignment and
\texttt{AWSIZE} are not checked at runtime.  \texttt{wr\_line\_count} $\le 64$
follows from the 4\,KB transaction limit (\S\ref{sec:overview}).  The
footprint is written at allocation (\S\ref{sec:alloc}), is never modified by
completion, and is released when the entry is freed (\S\ref{sec:release}).

\textbf{Read admission check.} At the head of the AR Metadata FIFO, on the
Request Validator valid path and \emph{before} \texttt{alloc\_req}, the
incoming read's footprint is computed from
\texttt{\{ARADDR, ARLEN, ARSIZE, ARBURST\}} by the identical calculation and
tested for interval overlap against every occupied \texttt{ROB\_WR} entry.
Writing $R$ for the incoming read footprint and $W$ for an occupied write
footprint, the read conflicts iff
\[
  (R_{\mathrm{base}} < W_{\mathrm{base}} + W_{\mathrm{count}})
  \;\wedge\;
  (W_{\mathrm{base}} < R_{\mathrm{base}} + R_{\mathrm{count}}).
\]
The read footprint is transient: used for the check only, stored nowhere.  No
footprint or hazard state is added to \texttt{ROB\_RD}, and a completed read
is never re-checked.

\textbf{Hold and release.} On any overlap the AR Metadata FIFO head is not
popped and no read \texttt{ROB\_INDEX} is allocated; the hold uses the same
mechanism as the existing per-AXID \texttt{STALL} (outputs held stable, FIFO
not advanced).  The check is re-evaluated whenever a write
\texttt{ROB\_INDEX} is freed (\texttt{free\_req}, \S\ref{sec:release}); the
freed entry's footprint stops being occupied in the same cycle its Free List
bit is set.  The read admits as soon as it overlaps no occupied write
footprint.  Liveness holds because writes drain independently of the held
read.  Holding the FIFO head also stalls non-conflicting reads behind it on
that port; this is accepted, as same-line overlap is rare.

\textbf{MC contract dependency.} Release keys on the blocking write's
\texttt{ROB\_WR} entry being freed, which occurs when its \texttt{BRESP} is
accepted after every \texttt{WR\_ACK} for that write has been received.  This
is correct only if a write reported complete is visible to a subsequent
same-line read; that property is stated as an explicit interface requirement
on MC Core (\S\ref{sec:cifmc-req}, item 2).

\section{CIF--MC Interface Contract}
\label{sec:cifmc}
\label{sec:pktformat}

This section is the authoritative statement of the CIF\,$\to$\,MC and
MC\,$\to$\,CIF interface.  Where an earlier section gives a placeholder
signal name or leaves an interface detail open, this section governs.  It
applies the frozen decisions of \S\ref{sec:identity}, \S\ref{sec:tag},
\S\ref{sec:rob}, and \S\ref{sec:samelinehaz} without changing them.

\subsection{Channels}

Three clock-domain-crossing FIFOs cross the CIF--MC boundary:

\begin{center}
\begin{tabular}{lll}
\toprule
\textbf{FIFO} & \textbf{Direction} & \textbf{Payload}\\
\midrule
\texttt{req\_afifo}   & CIF $\to$ MC & \texttt{\{req\_type[0], GLOBAL\_TAG[13:0], daddr[DADDR\_W-1:0]\}}\\
\texttt{wdata\_afifo} & CIF $\to$ MC & \texttt{\{wdata[511:0], wmask[63:0]\}} --- write packets only\\
\texttt{resp\_afifo}  & MC $\to$ CIF & \texttt{\{resp\_type[0], GLOBAL\_TAG[13:0], status[3:0], rdata[511:0]\}}\\
\bottomrule
\end{tabular}
\end{center}

\noindent \texttt{req\_afifo} is one ordered stream carrying read and write
descriptors interleaved in CIF packet-emission order; that order is program
order.  \texttt{wdata\_afifo} carries one entry per write descriptor, in the
same order the write descriptors enter \texttt{req\_afifo} (positional
binding); a write packet is emitted only when both request- and write-data
credit are available.  \texttt{resp\_afifo} is one merged stream;
\texttt{rdata} is valid only when \texttt{resp\_type = RD\_DATA}.

\subsection{Request packet}

\begin{center}
\begin{tabular}{lp{1.6cm}p{8.4cm}}
\toprule
\textbf{Field} & \textbf{Width} & \textbf{Meaning}\\
\midrule
\texttt{req\_type}   & 1 & \texttt{0 = READ}, \texttt{1 = WRITE}.  Invalid AXI requests never cross (see \S\ref{sec:cifmc} error transactions)\\
\texttt{GLOBAL\_TAG} & 14 & \texttt{\{ROB\_INDEX[7:0], PACKET\_NUMBER[5:0]\}} (\S\ref{sec:tag}); opaque to MC, echoed verbatim on the response\\
\texttt{daddr}       & \texttt{DADDR\_W} & DRAM coordinates \texttt{\{rank, ch, bg, bank, row, col\}}, MSB\,$\to$\,LSB, produced by the CIF-side address map/hash.  \emph{Not} a byte address; sub-block / offset bits excluded.  \texttt{DADDR\_W = RANK\_BITS + CH\_BITS + 3 + 2 + ROW\_BITS + COL\_BITS} (\texttt{BG\_BITS = 3}, \texttt{BANK\_BITS = 2} frozen); baseline 30\\
\bottomrule
\end{tabular}
\end{center}

\noindent One request packet is one 64-B AXI data beat; up to
\texttt{MAX\_PKTS = 64} per transaction.  \texttt{PACKET\_NUMBER} is the Burst
Splitter \texttt{pkt\_idx}, \texttt{0..exp\_pkts-1}.  There is no
\texttt{beat\_cnt}, \texttt{size}, or \texttt{len} field.  For WRAP bursts the
CIF-side WRAP Beat Address Table produces the per-packet byte address that is
then mapped to \texttt{daddr}.

\subsection{Write-data packet}

\begin{center}
\begin{tabular}{lp{1.6cm}p{8.4cm}}
\toprule
\textbf{Field} & \textbf{Width} & \textbf{Meaning}\\
\midrule
\texttt{wdata} & 512 & one 64-B write beat\\
\texttt{wmask} & 64  & byte-enable, bit \texttt{i} $\leftrightarrow$ byte \texttt{i}, verbatim from AXI \texttt{WSTRB}.  \texttt{wmask} not all-ones $\Rightarrow$ MC read-modify-write\\
\bottomrule
\end{tabular}
\end{center}

\subsection{Response packet}
\label{sec:respformat}

\begin{center}
\begin{tabular}{lp{1.6cm}p{8.4cm}}
\toprule
\textbf{Field} & \textbf{Width} & \textbf{Meaning}\\
\midrule
\texttt{resp\_type}  & 1  & \texttt{0 = RD\_DATA} (route to read ROB), \texttt{1 = WR\_ACK} (route to write ROB)\\
\texttt{GLOBAL\_TAG} & 14 & echoed verbatim; CIF decodes \texttt{ROB\_INDEX = [13:6]}, \texttt{PACKET\_NUMBER = [5:0]}\\
\texttt{status}      & 4  & \texttt{0 = OK}, \texttt{1 = SLVERR}, \texttt{2 = DECERR}.  Codes \texttt{3..15} reserved; CIF treats a reserved value as \texttt{DECERR}\\
\texttt{rdata}       & 512 & \texttt{RD\_DATA} only; undefined for \texttt{WR\_ACK}\\
\bottomrule
\end{tabular}
\end{center}

\noindent One MC completion sets exactly one \texttt{completion\_bitmap} bit,
one per 64-B packet.  A packet completing is not \texttt{RLAST}:
\texttt{RLAST} additionally requires \texttt{PACKET\_NUMBER == exp\_pkts - 1}
and that beat to have been serialized and accepted downstream.

\subsection{Flow control}

\textbf{Request side (\texttt{req\_afifo}, \texttt{wdata\_afifo}).}
Credit-based push, one independent credit pool per FIFO.  CIF holds
\texttt{N\_REQ\_CREDITS} (equal to the FIFO depth), drives a push while
credit $>$ 0, and decrements per push.  MC returns one \texttt{credit\_return}
pulse per consumed entry (registered, 2-flop synchronized MC\,$\to$\,CIF).
No combinational ready crosses the clock-domain boundary; there is no
\texttt{PKT\_READY}.

\textbf{Response side (\texttt{resp\_afifo}).}
Valid-only push; no ready from CIF.  MC self-throttles: every read reserves a
\texttt{resp\_afifo} slot at MC issue time and is not issued unless a slot is
free.  CIF accepts every response the cycle it is presented (one per cycle,
merged stream, dedicated Metadata-RAM completion port) and returns a response
credit to MC only after the response is fully processed --- the
\texttt{completion\_bitmap} bit set, and for \texttt{RD\_DATA} the
\texttt{rdata} written to \texttt{rd\_dbuf\_base + PACKET\_NUMBER}.

\subsection{Transaction identity}

The 14-bit \texttt{GLOBAL\_TAG} is the sole end-to-end identifier.  MC stores
it verbatim and never interprets any sub-field.  No AXI ID crosses the
interface; CIF recovers the AXID from per-AXID Slot-FIFO context at
retirement (\S\ref{sec:fifo}).

\subsection{Status to AXI response}

\begin{center}
\begin{tabular}{ll}
\toprule
\texttt{status} & AXI \texttt{RESP[1:0]}\\
\midrule
\texttt{0} & \texttt{OKAY} (\texttt{2'b00})\\
\texttt{1} & \texttt{SLVERR} (\texttt{2'b10})\\
\texttt{2} & \texttt{DECERR} (\texttt{2'b11})\\
\texttt{3..15} & reserved $\to$ \texttt{DECERR}\\
\bottomrule
\end{tabular}
\end{center}

\noindent \texttt{status} is per-packet.  \textbf{Read:} each packet's mapped
\texttt{RESP} drives \texttt{RRESP} for the AXI R beat it serializes into ---
no aggregation.  \textbf{Write:} \texttt{resp\_status} is the worst-of over
the transaction's \texttt{exp\_pkts} \texttt{WR\_ACK}s, precedence
\texttt{DECERR} $>$ \texttt{SLVERR} $>$ \texttt{OKAY} (numeric max of the
\texttt{status} codes); the single AXI \texttt{BRESP} is
\texttt{map(resp\_status)}.  \texttt{EXOKAY} is unused.

\subsection{Error transactions}

Invalid AXI requests (CIF address-range-check failures) are completed
entirely inside CIF and \emph{do not cross the interface}.  The Error Handler
drives \texttt{alloc\_req}; at Metadata-RAM initialization it force-fills
\texttt{completion\_bitmap[exp\_pkts-1:0]} and sets
\texttt{resp\_status = SLVERR}, so the entry is immediately
retirement-eligible and retires in Slot-FIFO order.  A read error emits
\texttt{exp\_pkts} R beats with \texttt{RRESP = SLVERR} and undefined data,
\texttt{rob\_last} on the final beat.  MC-detected faults on \emph{valid}
transactions are carried in the \texttt{status} field of the owning
\texttt{RD\_DATA}/\texttt{WR\_ACK}, not as a distinct response type.

\subsection{Requirements on MC Core}
\label{sec:cifmc-req}

The two items below are \textbf{requirements this interface places on MC
Core}.  They are stated as interface dependencies; the current MC-core
documentation is \emph{not} asserted to implement or to have verified them.

\begin{enumerate}
  \item \textbf{Response granularity.}  MC shall generate \emph{exactly one}
        response packet for every accepted 64-B request packet, echoing that
        packet's \texttt{GLOBAL\_TAG} unmodified.  No burst-level or
        transaction-level response grouping.  There is no \texttt{last} field
        on the interface --- CIF derives \texttt{RLAST} from
        \texttt{PACKET\_NUMBER == exp\_pkts - 1}.  Response packets for one
        transaction may arrive in any order; CIF's \texttt{completion\_bitmap}
        and \texttt{rd\_dbuf\_base + PACKET\_NUMBER} placement tolerate that.
  \item \textbf{Write-completion safety (Issue \#11 dependency).}  MC shall
        issue \texttt{WR\_ACK} for a write packet only when that packet is
        complete \emph{and} a subsequent read to the same 64-B address is
        guaranteed to observe the write.  CIF releases a held same-line read
        (\S\ref{sec:samelinehaz}) when the blocking write's
        \texttt{ROB\_INDEX} is freed --- which occurs once every
        \texttt{WR\_ACK} for that write has been accepted and its
        \texttt{BRESP} emitted.  The correctness of Issue \#11 therefore
        depends on this requirement; CIF adds no further MC-visibility
        guarantee of its own.
\end{enumerate}

\subsection{Ordering}

\texttt{req\_afifo} order is program order; MC derives its internal ordering
token from ingress order.  MC may schedule DRAM commands and return responses
out of order.  MC does not reorder the response stream for AXI purposes;
CIF owns AXI R- and B-channel ordering per AXID (\S\ref{sec:fifo}).

\section{Ordering, capacity, and reset}

Slot FIFO order is transaction acceptance order for each AXID and therefore
enforces AXI response ordering.  Metadata RAM capacity is transaction capacity
only.  A full Slot FIFO blocks allocation for that AXID; an empty Free List
blocks allocation globally.  The latter means all 256 transaction slots are
allocated, not merely that one Slot FIFO is full.  Read request issuance is
controlled separately by pooled Read Data Buffer availability.

Reset is active-low, asserted asynchronously, and deasserted synchronously;
\S\ref{sec:resetseq} gives the full initialization sequence.  It clears both
ROB Free Lists, Metadata RAM state, Slot FIFOs, and Retiring Registers.  The
Read Data Buffer is reset and drained according to its own specification.  A
soft reset is entered only after outstanding responses have been accepted and
both ROB instances have no allocated Slot FIFO entries.

\section{Error handling}

Invalid requests (CIF address-range-check failures) are completed entirely
inside CIF and never cross the CIF--MC interface (\S\ref{sec:cifmc}, error
transactions).  The Error Handler allocates the ROB entry, force-fills
\texttt{completion\_bitmap} and sets \texttt{resp\_status = SLVERR} at
initialization, and the entry retires in Slot-FIFO order.  The
\texttt{status} $\to$ AXI \texttt{RESP} mapping is frozen in \S\ref{sec:cifmc}
(\texttt{0}/\texttt{1}/\texttt{2} $=$ OK/SLVERR/DECERR); only the meaning of
MC error codes \texttt{3..15} remains an MC-owned detail.
